\section{Síntese do Trabalho}

Este trabalho apresentou o desenvolvimento e a avaliação experimental do WCAN, uma camada de comunicação sem fio baseada em ESP-NOW para transportar mensagens com semântica CAN entre módulos sensores e o sistema de aquisição de dados de um veículo de Fórmula SAE. A motivação principal foi reduzir a dependência de cabeamento em regiões de difícil instrumentação, como rodas, componentes aerodinâmicos e sensores externos ao veículo, sem exigir infraestrutura Wi-Fi, endereçamento IP ou mudanças profundas no modelo de dados já usado pela equipe.

A solução desenvolvida preserva o CAN ID como identificador lógico do sinal, agrupa amostras em pacotes ESP-NOW, usa filas FreeRTOS para desacoplar geração, transmissão e recepção, e inclui mecanismos de filtragem, confirmação e retransmissão. O trabalho também separou a camada WCAN da camada CAN física baseada no driver TWAI. Essa separação foi importante porque o driver CAN da Espressif já é uma base madura e validada, enquanto o ponto de incerteza deste projeto era a ponte sem fio. Os testes com transceptores VP230 em bancada indicaram que a integração CAN era funcional, mas a montagem em protoboard introduzia variabilidade suficiente para desviar a investigação do problema central. Por isso, a avaliação final concentrou-se no WCAN.

Os resultados da campanha \texttt{full\_run} mostram que o WCAN funciona bem em bancada quando usado no regime para o qual foi projetado: aquisição de dados com frequência conhecida, batching habilitado e tolerância a latências na ordem de milissegundos. A variante broadcast com \texttt{linger\_ms = 100} foi a configuração mais robusta. Ela manteve 100\% de entrega na suíte \texttt{baseline}, suportou múltiplas topologias de sensores e receptores e continuou operacional em frequências de até 1000 Hz na métrica principal de entrega.

O resultado mais importante da suíte \texttt{real\_time} foi mostrar que esse desempenho não vem simplesmente da taxa física do rádio. Sem batching, o sistema passa a ser limitado pela quantidade de pacotes ESP-NOW independentes, callbacks, filas, ACKs e retransmissões. Isso é coerente com o objetivo do projeto: o WCAN não pretende substituir uma rede CAN de controle em tempo real, mas expandir aquisição de dados em sensores amostrados com frequência conhecida. Nesse cenário, agrupar amostras por uma pequena janela de tempo é uma decisão aceitável e, pelos testes, necessária para operar em frequências altas.

A comparação entre variantes também deixa claro que a semântica de confiabilidade precisa ser escolhida de acordo com o uso. O broadcast é simples, robusto e adequado quando basta que ao menos um receptor interessado confirme o pacote, comportamento próximo à lógica de ACK dominante do CAN. Porém, essa semântica não garante que todos os receptores de uma topologia receberam todos os pacotes. Para esse requisito, a evolução natural é usar multicast com confirmação por assinante ou algum mecanismo equivalente de ACK por receptor.

\section{Problemas Conhecidos e Melhorias Futuras}

O principal problema conhecido da versão atual é o comportamento em situações sem assinantes válidos. Nos logs multicast, alguns sensores permaneceram tentando transmitir em modo de descoberta, mesmo quando nenhum assinante havia sido consolidado para aquele CAN ID. Em frequências altas, esse estado é perigoso: o sistema continua gerando pacotes, pressionando filas e heap, até registrar erros como \texttt{ESP\_ERR\_ESPNOW\_NO\_MEM} ou falhas de alocação dinâmica. Um sensor sem assinantes não deve derrubar a aplicação. O comportamento correto é tentar a descoberta por um número limitado de vezes, aplicar backoff, suspender temporariamente aquele fluxo ou descartar pacotes de forma controlada, registrar o evento e continuar executando.

Uma melhoria diretamente ligada a esse problema é reduzir alocações dinâmicas no caminho crítico de transmissão. A criação de pacotes por \texttt{new} ou \texttt{make\_unique} durante envio em alta frequência torna o sistema vulnerável a fragmentação, esgotamento de heap e falhas difíceis de reproduzir. Uma versão mais robusta deve usar buffers pré-alocados, pools de pacotes ou filas com capacidade fixa, de modo que o pior caso de memória seja conhecido antes da execução. Quando um pool estiver cheio, o sistema deve falhar de forma controlada, preferencialmente descartando a amostra mais nova ou mais antiga de acordo com a política definida, mas sem reiniciar a placa.

Também é necessário amadurecer a descoberta e a manutenção de assinantes. A transição entre broadcast de descoberta, registro de assinantes e envio multicast precisa ter estados explícitos, timeouts claros e métricas próprias. O \texttt{full\_run} mostrou que algumas falhas não eram perda de rádio em regime permanente, mas falhas de convergência: receptores podiam observar pacotes iniciais enquanto o transmissor não consolidava a tabela de assinantes. Uma melhoria futura é separar, nos logs e no protocolo, os eventos de descoberta, assinatura, expiração e reconexão, facilitando distinguir perda real de pacote de falha de controle.

Outra limitação é o dimensionamento de pilha e heap. O analisador encontrou high-water marks baixos em algumas tasks e, nos testes multicast, a ausência prolongada de assinantes válidos levou a pressão de memória suficiente para expor falhas de alocação. Mesmo quando esses casos não resultaram em crash, eles indicam pouca margem para uso embarcado prolongado. Antes de uma validação no veículo, é necessário revisar os tamanhos de stack, medir a ocupação por task em execuções longas, monitorar fragmentação de heap e transformar falhas de memória em erros recuperáveis.

A camada de análise também pode evoluir. A campanha \texttt{full\_run} evidenciou que uma repetição bem-sucedida não significa, sozinha, que a topologia é robusta, e uma repetição reprovada também não implica necessariamente falha sistemática do protocolo. O analisador já separa entrega, transporte, cenário, heap, crash e frequência observada, mas os resultados mostram a necessidade de classificar automaticamente falhas por causa provável: saturação, descoberta incompleta, receptor específico, ausência global de recepção, heap, pilha ou borda de teste. Isso tornaria a interpretação menos dependente de inspeção manual dos logs.

Do ponto de vista experimental, a próxima etapa é repetir as topologias críticas com mais repetições e duração maior. As suítes \texttt{mixed\_frequency} e algumas falhas isoladas do \texttt{baseline} apontam fragilidades, mas nem sempre têm repetições suficientes para separar variação experimental de erro sistemático. Testes longos também são importantes para validar se os alertas de heap são vazamentos reais, fragmentação acumulada ou apenas variação normal da execução.

Por fim, ainda falta a validação embarcada no veículo. A bancada é adequada para comparar variantes e encontrar gargalos de firmware, mas não reproduz vibração, alimentação real, posicionamento de antenas, interferência eletromagnética, massa final do encapsulamento e integração completa com o barramento CAN do carro. A partir dos resultados deste trabalho, a configuração mais indicada para essa etapa é broadcast com batching, seguida por uma nova rodada de multicast depois que o comportamento sem assinantes e o gerenciamento de memória forem corrigidos.

\section{Considerações Finais}

A contribuição principal deste trabalho foi demonstrar que uma ponte sem fio com semântica CAN é viável para aquisição de dados de Fórmula SAE quando seus limites são assumidos explicitamente. O WCAN não substitui o CAN físico em aplicações de controle crítico, mas pode reduzir cabeamento e permitir instrumentação modular em pontos onde a rede cabeada é difícil ou indesejável.

Os testes também foram úteis por mostrarem onde a solução ainda não está pronta. O broadcast com batching já é um candidato prático para uso controlado em bancada e futura validação no veículo. O multicast é conceitualmente o caminho certo para múltiplos receptores, mas precisa de uma política mais forte de descoberta, confirmação e memória. Com essas melhorias, o WCAN pode evoluir de protótipo funcional para uma camada de aquisição sem fio mais previsível, segura e integrada ao ecossistema de dados da equipe.
